% Part: lambda-calculus
% Chapter: lambda-definability
% Section: lists

\documentclass[../../../include/open-logic-section]{subfiles}

\begin{document}

\olfileid{lam}{ldf}{lst}
\olsection{لړونه}

لړ $[M_1, M_2, \ldots, M_n]$ په دې توګه تعريفولے شو:
\[
  [M_1, M_2, \ldots, M_n] \ident \lambd[sz][s M_1 (s M_2 (\ldots (s M_n z)))]
\]
په غير رسمي ژبه، د لړ تمثيلوونکے ترم د هغه لړ «راټولونکې» تابعه ده: داسې تابعه چې دوه ترمونه مني۔ لومړے يې داسې تابعه ده چې يو نوے عنصر او ورسره تر دې دمه راټوله شوې پايله اخلي، او نوې راټوله شوې پايله ورکوي۔ دويم يې د راټولېدنې لومړنے ارزښت دے۔ عناصر يو په يو راټولوي او په پاے کښې پايله ورکوي۔

\begin{digress}
که له تابعوي پروګرام جوړولو سره بلد ياست، نو دا تعريف به د \emph{ښي خوا راټولونې} سره ورته ومومئ: لړ د خپلو عناصرو د ښي خوا راټولونکې تابعې په توګه تعريفېږي۔
\end{digress}

د دې کوډونې په مرسته څو ګټورې تابعې په اسانۍ تعريفولے شو:
\begin{align*}
  \fn{Sum} & \ident
  \lambd[l][l \, (\lambd[xa][\fn{Add} \, x \, a])\,  \num{0}]\\
  \fn{Len} & \ident
  \lambd[l][l \, (\lambd[xa][\fn{Add} \, a \, \num{1}]) \num{0}]
\end{align*}

$\fn{Sum}$ د چرچ د عددنښو د يوه لړ مجموعه حسابوي۔ په لړ کښې راټولونه کوي: لومړنے ارزښت $\num{0}$ وي، او د هر عنصر دپاره $\fn{Add} \, x\, a$ د نوي ارزښت په توګه ورکوي؛ دلته $x$ اوسنے عنصر او $a$ تر دې دمه راټول شوے ارزښت دے۔ پايله د ټولو عناصرو مجموعه ده۔

\begin{prob}
  $\fn{Len}$ څۀ کوي؟ تشريح يې کړئ۔
\end{prob}

\begin{prob}
  داسې تابعه تعريف کړئ چې د يوه لړ ترتيب سرچپه کړي۔
\end{prob}
\end{document}
